\section{Where classical equivalences fail}\label{sec:notquantum}
Classical edge dropping and moderate face splitting can preserve $\C$ while shrinking $\Q$.
In each witness an outside party holds a quantum system correlated with the proposed replacement.
We write $a$ for the outcome of $A$, and similarly for the other parties.

\subsection{The classical edge rule}\label{sec:hlp-fails}

\Needspace{13\baselineskip}
\begin{lemma}[A party with a double is local]\label{lem:predict}
Let $Q(a,a',c\mid s,s',t)$ be a box for three parties $A$, $A'$, $C$, with inputs $s,s'\in\{0,1\}$ and $t$, that is
no-signalling: the joint distribution of the outcomes of any two parties does not depend on the input of the
third. Suppose that $A'$ behaves towards $C$ exactly as $A$ does: for $i\in\{0,1\}$ and all $\alpha,c,t$,
\[Q(a'=\alpha,\,c\mid s'=i,\,t)\;=\;Q(a=\alpha,\,c\mid s=i,\,t)\;=:\;P(\alpha,c\mid i,t).\]
Then $P$ is local: there are a distribution $\nu(a_0,a_1)$ and response functions $P(c\mid t,a_0,a_1)$ with
\[
P(a,c\mid s,t)=\sum_{a_0,a_1}\nu(a_0,a_1)\,\delta_{a,a_s}\,P(c\mid t,a_0,a_1).
\]
In particular $P$ satisfies the CHSH inequality.
\end{lemma}
\begin{proof}
Put $\mu_t(a_0,a_1,c)=Q(a=a_0,\,a'=a_1,\,c\mid s=0,\,s'=1,\,t)$. Its marginal over $a_1$ is $P(a_0,c\mid0,t)$, by
no-signalling from $A'$. Its marginal over $a_0$ is $Q(a'=a_1,c\mid s'=1,t)=P(a_1,c\mid1,t)$, by no-signalling
from $A$ and the hypothesis. Its marginal over $c$ is a distribution $\nu(a_0,a_1)$ that does not depend on $t$,
by no-signalling from $C$. So $\lambda=(a_0,a_1)$, drawn from $\nu$, with $A$ answering $a_s$ and $C$ answering
according to $\mu_t(\lambda,c)/\nu(\lambda)$ (arbitrary where $\nu(\lambda)=0$), reproduces $P$.
\end{proof}

\begin{cexbox}{the edge rule, with settings}
\begin{theorem}\label{thm:hlp-settings}
Let $\G^{+}$ have the nodes $A,B,C,S,T$, the edges $S\to A$, $S\to B$, $T\to C$, $A\to B$ and the source
$\{A,B,C\}$, and let $\G^{-}=\G^{+}-(A\to B)$. Then
\[
\C(\G^{-})=\C(\G^{+})\qquad\text{but}\qquad\Q(\G^{-})\subsetneq\Q(\G^{+}).
\]
\end{theorem}
\end{cexbox}

\begin{figure}[H]
\centering
\begin{tikzpicture}[mdag]
  \begin{scope}
    \vn{s}{0.9,1.9}{$S$} \vn{a}{0,0.6}{$A$} \vn{b}{1.8,0.6}{$B$} \vn{c}{0.9,-1.0}{$C$} \vn{t}{2.3,-1.0}{$T$}
    \qs{q}{0.9,-0.05}
    \feed{q}{a,b,c} \de{s}{a} \de{s}{b} \de{t}{c} \de{a}{b}
    \node[tag] at (0.9,-1.75) {$\G^{+}$: $B$ copies $A$};
  \end{scope}
  \node[rel] at (3.5,0.4) {$\ceq$};
  \node[rel] at (3.5,-0.25) {$\not\qeq$};
  \begin{scope}[shift={(4.6,0)}]
    \vn{s}{0.9,1.9}{$S$} \vn{a}{0,0.6}{$A$} \vn{b}{1.8,0.6}{$B$} \vn{c}{0.9,-1.0}{$C$} \vn{t}{2.3,-1.0}{$T$}
    \qs{q}{0.9,-0.05}
    \feed{q}{a,b,c} \de{s}{a} \de{s}{b} \de{t}{c}
    \node[tag] at (0.9,-1.75) {$\G^{-}$: $B$ would have to guess};
  \end{scope}
\end{tikzpicture}
\caption{With $A\to B$, $B$ can copy Alice while Alice and Charlie play CHSH. Without the edge, a successful copy would supply a second Alice.}
\label{fig:hlp-settings}
\end{figure}

\begin{proof}[Proof of \cref{thm:hlp-settings}]
(H1) and (H2) give classical equality. With the edge, give $A,C$ a Bell pair, play the Clauser--Horne--Shimony--Holt (CHSH) game at independent
uniform settings $s,t$, and let $B$ copy $a$. Without the edge, $A,B,C$ would perform local measurements
on a shared state. Give $B$ its own input $s'$ to select the same measurement it would perform at $s=s'$.
The resulting box is no-signalling. When $s=s'$ we have $b=a$, so by no-signalling the $BC$ marginal
at $s'=i$ equals the $AC$ marginal at $s=i$ for either $i$. \Cref{lem:predict} makes $AC$ local,
contradicting the observed CHSH value $2\sqrt2$. Structural dominance gives the other inclusion.
\end{proof}

\begin{cexbox}{the edge rule, on four nodes without settings}
\begin{theorem}\label{thm:hlp-four}
Let $\G^{+}$ have the nodes $A,B,C,S$, the sources $\{A,B,C\}$, $\{A,B,S\}$, $\{C,S\}$ and the single edge
$A\to B$, and let $\G^{-}$ be the same without the edge. Then $\C(\G^{-})=\C(\G^{+})$ but
$\Q(\G^{-})\subsetneq\Q(\G^{+})$.
\end{theorem}
\end{cexbox}

\emph{The distribution.} $P^{*}$ is Fritz's distribution for the triangle \cite{Fritz}, with $B$ copying $A$. The
source $\{A,B,C\}$ gives $A$ and $C$ a Bell pair; $\{A,B,S\}$ gives $A$ and $S$ a uniform bit $x$; $\{C,S\}$ gives
$C$ and $S$ a uniform bit $y$. $A$ measures with the CHSH setting $x$ and outputs $(a,x)$; $C$ measures with the
setting $y$ and outputs $(c,y)$; $S$ outputs $(x,y)$; $B$ outputs what $A$ outputs. Then
$P^{*}(a,x,c,y)=\tfrac14P_{\mathrm{CHSH}}(a,c\mid x,y)$ with CHSH value $2\sqrt2$, and $P^{*}\in\Q(\G^{+})$.

\begin{figure}[H]
\centering
\begin{tikzpicture}[mdag]
  \begin{scope}[shift={(0,0.1)}]
    \vn{a}{0,0}{$A$} \vn{b}{2.8,0}{$B$} \vn{c}{1.4,1.9}{$C$} \vn{s}{1.4,-1.9}{$S$}
    \qs{f1}{1.4,0.85} \qs{f2}{1.4,-0.85} \qs{f3}{3.9,0}
    \feed{f1}{a,b,c} \feed{f2}{a,b,s} \feed{f3}{c,s}
    \node[lab,text=qcol] at (1.4,0.42) {$\rho_1$}; \node[lab,text=qcol] at (1.4,-0.42) {$\rho_2$};
    \node[lab,text=qcol] at (4.35,0) {$\rho_3$};
    \node[tag] at (1.4,-2.75) {$\G^{-}$};
  \end{scope}
  \begin{scope}[shift={(7.4,0.3)}]
    % sources of the experiment
    \qs{r1}{2.4,1.0}
    \qs{r2}{0.3,1.9} \qs{r2p}{4.5,1.9} \qs[src,faint]{r2pp}{6.6,-0.2}
    \qs{r3}{2.4,-1.3} \qs[src,faint]{r3p}{-1.6,0.6} \qs[src,faint]{r3pp}{6.4,1.5}
    \vn{A}{1.1,1.0}{$\bar A$} \vn{B}{3.7,1.0}{$\bar B$} \vn{C}{2.4,-0.2}{$\bar C$}
    \vn{SA}{-0.6,1.0}{$\bar S_{\!A}$} \vn{SB}{5.5,1.0}{$\bar S_{\!B}$} \vn{SC}{4.4,-1.0}{$\bar S_{\!C}$}
    \feed{r1}{A,B,C} \feed{r2}{A,SA} \feed{r2p}{B,SB} \feed{r3}{C,SC}
    \draw[qedge,faint] (r3p) -- (SA); \draw[qedge,faint] (r3pp) -- (SB); \draw[qedge,faint] (r2pp) -- (SC);
    \node[lab,text=qcol] at (2.4,1.4) {$\rho_1$}; \node[lab,text=qcol] at (0.3,2.3) {$\rho_2$};
    \node[lab,text=qcol] at (4.5,2.3) {$\rho_2'$}; \node[lab,text=qcol] at (2.4,-1.7) {$\rho_3$};
    \node[lab,text=qcol!60] at (-1.6,0.2) {$\rho_3'$}; \node[lab,text=qcol!60] at (6.4,1.9) {$\rho_3''$};
    \node[lab,text=qcol!60] at (6.6,-0.6) {$\rho_2''$};
    \node[tag] at (2.4,-2.4) {the experiment: $\bar A$ and $\bar B$ get their bits from different copies};
  \end{scope}
\end{tikzpicture}
\caption{The four-node witness and its inflation. Giving $\bar B$ a fresh copy of the setting source
preserves its correlations with $C$ while separating its setting from $A$'s.}
\label{fig:inflation}
\end{figure}

\begin{proof}
\emph{Classical.} This is the classical edge rule again: $A$ has no setting, and both its sources contain $B$.

\emph{Quantum.} $P^{*}\in\Q(\G^{+})$ as described. Suppose a quantum model of $\G^{-}$ produces $P^{*}$: states
$\rho_1$ on the shares $A_1B_1C_1$, $\rho_2$ on $A_2B_2S_2$, $\rho_3$ on $C_3S_3$, and POVMs $M_A$ (on $A_1A_2$),
$M_B$ (on $B_1B_2$), $M_C$ (on $C_1C_3$), $M_S$ (on $S_2S_3$). Consider the following experiment, of the kind
used in the inflation technique \cite{Inflation}, built from $\rho_1$, three independent copies
$\rho_2,\rho_2',\rho_2''$ of the second state and three copies $\rho_3,\rho_3',\rho_3''$ of the third:
\begin{center}\small
\begin{tabular}{@{}lll@{}}
\toprule
party & measurement, on the shares & outcome\\
\midrule
$\bar A$ & $M_A$ on $A_1$ of $\rho_1$ and $A_2$ of $\rho_2$ & $(a,x)$\\
$\bar C$ & $M_C$ on $C_1$ of $\rho_1$ and $C_3$ of $\rho_3$ & $(c,y)$\\
$\bar B$ & $M_B$ on $B_1$ of $\rho_1$ and $B_2$ of $\rho_2'$ & $(a',x')$\\
$\bar S_A$ & $M_S$ on $S_2$ of $\rho_2$ and $S_3$ of $\rho_3'$ & first component $u$\\
$\bar S_B$ & $M_S$ on $S_2$ of $\rho_2'$ and $S_3$ of $\rho_3''$ & first component $u'$\\
$\bar S_C$ & $M_S$ on $S_2$ of $\rho_2''$ and $S_3$ of $\rho_3$ & second component $w$\\
\bottomrule
\end{tabular}
\end{center}
Each share is measured at most once, so the six parties have a joint distribution $P_{\rm exp}$.
The following marginals use the same reduced states as their counterparts in
$P^*$. Thus $(a,x,c,y)$ and $(a',x',c,y)$ both have the CHSH distribution with uniform settings,
and $x=u$, $x'=u'$, $y=w$ almost surely. Each auxiliary party $\bar S_i$ uses sources disjoint from
the other two main parties. Substituting its perfectly copied bit therefore gives
\[
 x'\perp(a,x,c,y),\qquad x\perp(a',x',c,y),\qquad y\perp(a,x,a',x').
\]
The three inputs are independent uniform bits, and these independences give no-signalling for the
conditional box $P_{\rm exp}(a,a',c\mid x,x',y)$. Its two CHSH marginals coincide, contradicting
\cref{lem:predict}. Thus one classical class on four nodes contains two different quantum model sets.
\end{proof}

\subsection{Moderate and strong face splitting}\label{sec:leaky}

\begin{cexbox}{moderate face splitting, on six nodes}
\begin{theorem}\label{thm:sixnode}
Let $\G_6$ have the nodes $A,R,D,E$ and the setting nodes $S,T$, the edges
\[S\to A,\quad S\to D,\quad T\to E,\quad R\to D,\quad A\to D,\]
and the sources $F=\{A,R,D\}$ and $J=\{R,D,E\}$. Let $\G_6'$ be the same with $F$ replaced by $\{A,R\}$. Then
\[
\C(\G_6')=\C(\G_6)\qquad\text{but}\qquad\Q(\G_6')\subsetneq\Q(\G_6).
\]
\end{theorem}
\end{cexbox}

\begin{figure}[H]
\centering
\begin{tikzpicture}[mdag]
  \begin{scope}
    \vn{s}{-0.4,3.0}{$S$} \vn[Cn]{A}{0,1.7}{$A$} \vn[Cn]{R}{2.2,1.7}{$R$} \vn[Dn]{D}{1.1,0}{$D$} \vn{E}{3.7,0}{$E$}
    \vn{t}{4.9,1.1}{$T$}
    \qs[splitsrc]{f}{1.1,1.15} \qs{g}{2.4,0.55}
    \feed[splitedge]{f}{A,R,D} \feed{g}{R,D,E}
    \de{A}{D} \de{R}{D} \de{s}{A} \de{t}{E}
    \draw[edge] (s) to[out=-150,in=165,looseness=1.3] (D);
    \node[lab,text=splitcol!80!black] at (1.1,1.62) {$F$}; \node[lab,text=qcol] at (3.0,0.85) {$J$};
    \node[tag] at (1.8,-0.75) {$\G_6$};
  \end{scope}
  \node[rel] at (6.0,1.3) {$\ceq$};
  \node[rel] at (6.0,0.6) {$\not\qeq$};
  \begin{scope}[shift={(7.9,0)}]
    \vn{s}{-0.4,3.0}{$S$} \vn[Cn]{A}{0,1.7}{$A$} \vn[Cn]{R}{2.2,1.7}{$R$} \vn[Dn]{D}{1.1,0}{$D$} \vn{E}{3.7,0}{$E$}
    \vn{t}{4.9,1.1}{$T$}
    \qs{f}{1.1,2.45} \qs{g}{2.4,0.55}
    \feed{f}{A,R} \feed{g}{R,D,E}
    \de{A}{D} \de{R}{D} \de{s}{A} \de{t}{E}
    \draw[edge] (s) to[out=-150,in=165,looseness=1.3] (D);
    \node[lab,text=qcol] at (3.0,0.85) {$J$};
    \node[tag] at (1.8,-0.75) {$\G_6'$};
  \end{scope}
\end{tikzpicture}
\caption{In both structures the relay belongs to $J$. Splitting $F$ removes Dan from Alice's source, even though Dan still sees Alice, her setting and the relay.}
\label{fig:sixnode}
\end{figure}

Dan ($D$) sees Alice ($A$), the relay ($R$) and Alice's setting $S$. Classically he can sample the latent of $F$ from its posterior
given these data and the latent of $J$ \cite{AWS}. The relay belongs to both sources and can measure its two shares jointly.
In the unsplit structure a relay measurement followed by Alice's measurement can leave Dan entangled with Eve ($E$), outside $F$.

\begin{figure}[H]
\centering
\begin{tikzpicture}[mdag,scale=0.82]
  \tikzset{link/.style={decorate, decoration={snake, amplitude=1.1pt, segment length=5pt}, line width=0.9pt},
           party/.style={draw=black!70, rounded corners=2pt, fill=black!4}}
  \newcommand{\parties}{
    \node[party, minimum width=7mm, minimum height=14mm] at (0,0.8) {}; \node[lab] at (0,1.9) {Alice};
    \coordinate (a1) at (0,1.15); \coordinate (a2) at (0,0.45);
    \node[party, minimum width=13mm, minimum height=6mm] at (1.7,2.3) {}; \node[lab] at (1.7,2.95) {relay};
    \coordinate (r1) at (1.4,2.3); \coordinate (r2) at (2.0,2.3);
    \node[party, minimum width=6mm, minimum height=6mm] at (3.3,1.15) {}; \node[lab] at (3.3,1.8) {Eve};
    \coordinate (e) at (3.3,1.15);
    \node[party, minimum width=6mm, minimum height=6mm] at (1.7,-0.3) {}; \node[lab] at (1.7,-0.95) {Dan};
    \coordinate (d) at (1.7,-0.3);
    \foreach \q in {a1,a2,r1,r2,e,d} { \fill (\q) circle (1.7pt); }
    \node[font=\tiny,anchor=east] at (-0.08,1.15) {1}; \node[font=\tiny,anchor=east] at (-0.08,0.45) {2};}
  \begin{scope}
    \parties
    \draw[link, splitcol!90!black] (a1) -- (r1); \draw[link, splitcol!90!black] (a2) -- (d); \draw[link, qcol] (r2) -- (e);
    \node[tag,align=center] at (1.7,-1.85) {three Bell pairs:\\ two from $F$, one from $J$};
  \end{scope}
  \begin{scope}[shift={(4.8,0)}]
    \parties
    \draw[link] (a1) -- (e); \draw[link, splitcol!90!black] (a2) -- (d);
    \node[lab] at (1.7,2.3) {$r$};
    \node[tag,align=center] at (1.7,-1.85) {the relay's Bell\\ measurement, outcome $r$};
  \end{scope}
  \begin{scope}[shift={(9.6,0)}]
    \parties
    \draw[link] (a1) -- (e); \draw[link, splitcol!90!black, faint] (a2) -- (d);
    \node[tag,align=center] at (1.7,-1.85) {$s\in\{0,1\}$: CHSH test\\ of Alice and Eve};
  \end{scope}
  \begin{scope}[shift={(14.4,0)}]
    \parties
    \draw[link] (d) -- (e);
    \node[lab] at (0,0.8) {$a$};
    \node[tag,align=center] at (1.7,-1.85) {$s=2$: Alice's Bell\\ measurement, outcome $a$};
  \end{scope}
\end{tikzpicture}
\caption{The strategy that produces $P^{*}$ in $\G_6$. Wavy lines are Bell pairs; black dots are qubits. The
relay's measurement links Alice to Eve. Alice then either tests CHSH against Eve, or passes the link on to Dan,
for which she needs the pair that $F$ gave her and Dan.}
\label{fig:strategy}
\end{figure}

\paragraph{A witness with two uses for one Bell pair.}
Take independent uniform inputs $s\in\{0,1,2\}$ and $t\in\{0,1\}$. Source $F$ supplies Alice--relay
and Alice--Dan Bell pairs, and $J$ supplies a relay--Eve pair. The relay's Bell measurement swaps
entanglement to Alice and Eve and announces the Bell label $r$. Eve measures $Z$ or $X$ at $t=0,1$.
At $s=0,1$, Alice measures $(Z+X)/\sqrt2$ or $(Z-X)/\sqrt2$ on her first qubit, recording
$0$ for $+1$ and $1$ for $-1$; Dan outputs $+1$. At $s=2$ Alice Bell-measures her two qubits, swapping
entanglement to Dan and Eve. Knowing $r,a$, Dan corrects the sign of his $Z$ outcome.

\Needspace{10\baselineskip}
The proof uses two observable properties, checked numerically in \cref{app:code}.
\begin{description}[leftmargin=1.2cm,style=nextline,itemsep=1pt]
\item[(P1)] $r$ is uniform, and for each $r$ Alice ($s\in\{0,1\}$) and Eve reach the CHSH value $2\sqrt2$, with a
  different pattern of signs for each $r$. Precisely, give Alice the sign $+1$ when $a=0$ and $-1$ otherwise, and
  let $c_r(s,t)$ be the expectation of the product of her sign and Eve's outcome, given $r$, $s$ and $t$. There are
  signs $\varepsilon(r)\in\{\pm1\}^2$, a different pair for each $r$, with
  \[\varepsilon_0c_r(0,0)+\varepsilon_1c_r(0,1)+\varepsilon_0c_r(1,0)-\varepsilon_1c_r(1,1)=2\sqrt2 .\]
\item[(P2)] at $s=2$, $t=0$: $d=e$ with probability one, and $e$ is uniform given $(r,a)$.
\end{description}

In particular, at $s=2$ Dan predicts a bit that is still uniform to anyone knowing only $(r,a)$.
We use the following exact form of CHSH rigidity \cite{SW}, proved in \cref{app:rigidity}.

\begin{lemma}[CHSH rigidity]\label{lem:rigidity}
Let $A_0,A_1$ be Hermitian operators on $\Hil_A$ and $E_0,E_1$ Hermitian operators on $\Hil_E$, all with square
$\id$, and let $\Hil_{\mathrm{ext}}$ be a further system; all spaces are finite-dimensional. For
$\varepsilon\in\{\pm1\}^2$ put
$W_\varepsilon=A_0\otimes(\varepsilon_0E_0+\varepsilon_1E_1)+A_1\otimes(\varepsilon_0E_0-\varepsilon_1E_1)$, and
suppose that for each $\varepsilon$ there is a state $\omega_\varepsilon$ on
$\Hil_A\otimes\Hil_E\otimes\Hil_{\mathrm{ext}}$ with
$\Tr[W_\varepsilon\,\omega_\varepsilon]=2\sqrt2$. Let $V_A\subseteq\Hil_A$ and $V_E\subseteq\Hil_E$ be the spans of
the supports of the reduced states of the $\omega_\varepsilon$. Then $V_A\cong\mathbb{C}^2\otimes\Hil_{A''}$ and
$V_E\cong\mathbb{C}^2\otimes\Hil_{E''}$ in such a way that on these subspaces
\[
A_0=\tfrac{Z+X}{\sqrt2}\otimes\id,\quad A_1=\tfrac{Z-X}{\sqrt2}\otimes\id,\quad E_0=Z\otimes\id,\quad E_1=X\otimes\id,
\qquad\text{and}\qquad \omega_\varepsilon=\proj{\beta_\varepsilon}\otimes\gamma_\varepsilon ,
\]
where $\beta_\varepsilon$ is the Bell state of two qubits with $Z\otimes Z=\varepsilon_0$ and
$X\otimes X=\varepsilon_1$, and $\gamma_\varepsilon$ is a state on
$\Hil_{A''}\otimes\Hil_{E''}\otimes\Hil_{\mathrm{ext}}$.
\end{lemma}

\begin{proof}[Proof of \cref{thm:sixnode}]
Classical equality follows from the moderate rule, and the witness above belongs to $\Q(\G_6)$.
Suppose a model of $\G_6'$ reproduces it. The independent sources have states $\rho_{AR_1}$ and
$\rho_{R_2DE}$. Alice's first two binary measurements (coarse-graining $a\ne0$) and Eve's two
measurements may be taken projective: use a private pure ancilla for each measurement, including it
in that party's source. This preserves source independence and leaves Alice's $s=2$ measurement available.

\emph{Average over the relay.} Its measurement on $R_1R_2$ leaves a conditional state $\omega_r$ on
$ADE$. Completeness of its POVM and independence of the sources give
$\sum_r p_r\omega_r=\rho_A\otimes\rho_{DE}$. By (P1), $p_r=1/4$ and the four conditioned states
attain the four signed CHSH maxima with the same observables.

\emph{Resolve the four Bell labels.} By \cref{lem:rigidity} the same qubit factors $A',E'$ work for
all $r$ on the supports of the averaged marginals. Eve's observables act as $Z$ and $X$ on $E'$ alone.
The conditional states are $\omega_r=\proj{\beta_r}_{A'E'}\otimes\gamma_r$ on those supports, hence
\[
 \tfrac14\sum_r\proj{\beta_r}_{A'E'}\otimes\gamma_r=\rho_A\otimes\rho_{DE}.
\]
Both Bell marginals are maximally mixed. Taking marginals of this identity gives
$\rho_A=\id_{A'}/2\otimes\alpha_{A''}$ and $\rho_{DE}=\id_{E'}/2\otimes\eta_{E''D}$.
Projecting onto each of the four orthogonal Bell states now gives $\gamma_r=\alpha\otimes\eta$
for every $r$.

\emph{Test Dan's prediction.} At $s=2$ Alice can measure both $A'$ and $A''$. Conditional on $(r,a)$,
this can steer $E'$, but leaves it in product with $\eta_{E''D}$. Eve measures $Z$ on $E'$ at $t=0$,
so her outcome is independent of Dan's given $(r,a)$. By (P2) that outcome is uniform, making their
agreement probability $1/2$ instead of one. This excludes the witness from $\Q(\G_6')$.
\end{proof}

The four-valued relay matters here. It identifies each Bell state uniquely, which is why the projection
fixes every $\gamma_r$. With sixteen outcomes the relay could swap two pairs and report two results.
Several labels would then give the same CHSH signs, and the projection would fix only their sum.

\paragraph{Four-node splits.}\label{sec:open-four}
On four nodes there are five groups for which the classical strong rule licenses a split and the
quantum rules give no such equivalence. The six-node witness uses an independent setting on each side and four different relay outcomes.
The four-node groups leave no room for the extra setting nodes. In the last two groups every node of
$C$ also belongs to a remaining source. Whether any of these splits preserves $\Q$ is open.

\begin{figure}[!htb]
\centering
\begin{tikzpicture}[mdag,scale=0.85]
  % group 1
  \begin{scope}
    \vn[Cn]{n0}{1.2,1.9}{$0$} \vn{n1}{3.1,1.9}{$1$} \vn[Cn]{n2}{-0.2,0.3}{$2$} \vn[Dn]{n3}{2.3,-0.1}{$3$}
    \qs[splitsrc]{f}{0.95,0.75} \qs{g}{2.4,1.2}
    \feed[splitedge]{f}{n0,n2,n3} \feed{g}{n0,n1,n3}
    \de{n0}{n1} \de{n0}{n2} \de{n2}{n3} \draw[edge] (n0) to[bend left=10] (n3);
    \node[tag] at (1.4,-0.9) {1};
  \end{scope}
  % groups 2 and 3
  \foreach \shift/\name in {5.6/2, 11.0/3} {
  \begin{scope}[shift={(\shift,0)}]
    \vn[Cn]{n0}{1.2,1.9}{$0$} \vn[Cn]{n1}{-0.7,0.95}{$1$} \vn[Dn]{n2}{1.2,-0.1}{$2$} \vn{n3}{3.1,-0.1}{$3$}
    \qs[splitsrc]{f}{0.5,0.92} \qs{g}{2.05,0.75}
    \feed[splitedge]{f}{n0,n1,n2} \feed{g}{n0,n2,n3}
    \de{n0}{n1} \de{n1}{n2} \de{n2}{n3} \draw[edge] (n0) to[bend left=12] (n2);
    \node[tag] at (1.2,-0.9) {\name};
  \end{scope}}
  \draw[edge] (n0) to[bend left=30] (n3);   % group 3 only: the nodes of the last panel drawn
  % groups 4 and 5: the green node is a member of all three sources
  \begin{scope}[shift={(1.2,-4.6)}]
    \vn[Cn]{n1}{0,1.2}{$1$} \vn[Cn]{n2}{4,1.2}{$2$} \vn[Dn]{n3}{2,1.2}{$3$} \vn{n0}{2,-0.9}{$0$}
    \qs[splitsrc]{f}{2,2.4} \qs{g1}{1.0,0.2} \qs{g2}{3.0,0.2}
    \feed[splitedge]{f}{n1,n2,n3} \feed{g1}{n0,n1,n3} \feed{g2}{n0,n2,n3}
    \de{n1}{n3} \de{n2}{n3}
    \node[tag] at (2,-1.7) {4};
  \end{scope}
  \begin{scope}[shift={(8.4,-4.6)}]
    \vn[Cn]{n0}{0,1.2}{$0$} \vn[Cn]{n1}{4,1.2}{$1$} \vn[Dn]{n2}{2,1.2}{$2$} \vn{n3}{2,-0.9}{$3$}
    \qs[splitsrc]{f}{2,2.4} \qs{g1}{1.0,0.2} \qs{g2}{3.0,0.2}
    \feed[splitedge]{f}{n0,n1,n2} \feed{g1}{n0,n2,n3} \feed{g2}{n1,n2,n3}
    \de{n0}{n2} \de{n1}{n2} \de{n2}{n3}
    \node[tag] at (2,-1.7) {5};
  \end{scope}
\end{tikzpicture}
\caption{The largest structure of each open group. The strong rule splits the orange source into the blue and green parts. In each graph a blue node also belongs to a remaining source.}
\label{fig:open-four}
\end{figure}
